\section*{ВВЕДЕНИЕ}
\addcontentsline{toc}{section}{ВВЕДЕНИЕ}

Учебно-методические центры повышения квалификации и профессиональной переподготовки играют ключевую роль в обеспечении профессионального развития сотрудников различных отраслей. В условиях стремительных изменений в технологиях и рыночных реалиях необходимость в непрерывном обучении становится особенно важной. Однако традиционный подход к организации учебного процесса, основанный на разрозненных электронных таблицах и ручном ведении отчётности, исчерпал свои возможности. Отсутствие централизованного хранилища данных, трудоёмкость формирования статистических сводок, сложность контроля за успеваемостью слушателей и отсутствие разграничения прав доступа существенно снижают эффективность работы центра.

Автоматизация процессов управления учебными курсами позволяет централизовать хранение данных, устранить ошибки ручного ввода, обеспечить оперативный контроль и существенно сократить время на подготовку отчётности. Проведённый анализ существующих систем (Moodle, Google Classroom, 1С:Университет) показал, что они либо обладают избыточным функционалом, либо не предоставляют специализированных средств для детального учёта контингента слушателей и формирования утверждённых форм статистической отчётности. В связи с этим целесообразной является разработка собственной специализированной информационной системы.

\emph{Цель настоящей работы} – разработка программно-информационной системы для автоматизации деятельности учебно-методического центра повышения квалификации. Для достижения поставленной цели необходимо решить \emph{следующие задачи:}
\begin{itemize}
	\item провести анализ предметной области и обосновать необходимость автоматизации;
	\item разработать концептуальную модель и спроектировать базу данных;
	\item спроектировать архитектуру программной системы (клиент-сервер с веб-интерфейсом и десктоп-приложением);
	\item реализовать систему средствами выбранных технологий (Python, Flask, PostgreSQL, Tkinter);
	\item выполнить системное тестирование разработанного приложения.
\end{itemize}

\emph{Структура и объем работы.} Отчет состоит из введения, 4 разделов основной части, заключения, списка использованных источников, 2 приложений. Текст выпускной квалификационной работы равен \formbytotal{lastpage}{страниц}{е}{ам}{ам}.

\emph{Во введении} сформулирована цель работы, поставлены задачи разработки, описана структура работы, приведено краткое содержание каждого из разделов.

\emph{В первом разделе} на стадии анализа предметной области приводится характеристика деятельности учебно-методического центра, выявляются проблемы традиционного подхода, обосновывается необходимость автоматизации и анализируются существующие аналоги.

\emph{Во втором разделе} на стадии технического задания формулируются требования к функциональности, интерфейсу, надёжности, информационной безопасности и совместимости, а также разрабатываются варианты использования (UML-диаграммы прецедентов).

\emph{В третьем разделе} на стадии технического проектирования представлены архитектура программной системы, обоснование выбора технологий (Python, Flask, PostgreSQL, Tkinter), проект базы данных (ER-модель, реляционная модель, структура таблиц).

\emph{В четвертом разделе} приводится спецификация классов и модулей разработанного приложения, а также результаты системного тестирования в виде формальных тест-кейсов.

В заключении излагаются основные результаты работы, полученные в ходе разработки.

В приложении А представлен графический материал.
В приложении Б представлены фрагменты исходного кода.